Die regulären Punkte erhält man durch das totale Differential: Die singulären Punkte bilden einen Kreis mit dem Radius R in der x-y-Ebene mit dem Nullpunkt in der Mitte (der der einzie singuläre Punkt ist, der nicht auf dem Kreis liegt), alle anderen Punkte sind regulär.

Die nebenstehende Animation zeigt beispielhaft für R=2 die Fasern für verschiedene s. Deutlich zu sehen ist, wie sich bei s=4=√(R) das Loch im Torus schließt.
[[File:Aufgabe79.27.gif|thumb|Animation zur Veranschaulichung von
[[Torus/Faserrealisierung/Andere Fasern/Aufgabe|Aufgabe]].
]]
Von Lukas Freudenberg. [[Kategorie:Latexseite]]
